\documentclass[11pt,a4paper]{article}

\usepackage[utf8]{inputenc}
\usepackage[T1]{fontenc}
\usepackage[italian]{babel}
\usepackage{amsmath,amssymb}
\usepackage{booktabs}
\usepackage{array}
\usepackage[margin=2.5cm]{geometry}

\setcounter{secnumdepth}{1}

\newcommand{\file}[1]{\texttt{\detokenize{#1}}}
\newcommand{\Tc}{T_c}
\newcommand{\nom}{\mathrm{nom}}
\newcommand{\lib}{\mathrm{lib}}
\newcommand{\stab}{s^{\mathrm{tab}}}        % punti di SOC delle tabelle
\newcommand{\rif}{[L]}                      % Linearizzazione_MPC.pdf
\newcommand{\rifq}{[Q]}                     % Problema_QP_MPC.pdf
\newcommand{\ip}[1]{\textbf{(H#1)}}         % etichetta di un'ipotesi

\title{Dal modello della cella al problema quadratico:\\passaggi e ipotesi}
\author{MPC di bilanciamento attivo, modello \file{Bilanciamento_MPC.slx}}
\date{}

\begin{document}
\maketitle

\noindent
Questo documento elenca, in ordine, i passaggi matematici che portano dal
modello della cella al problema quadratico (QP) che l'MPC di bilanciamento
risolve a ogni passo di controllo. Per ogni passaggio riporta la formula e le
ipotesi su cui si regge. Dimostrazioni ed esempi numerici sono nei due
documenti di riferimento:
\begin{itemize}
  \item {\rif} $=$ \file{Linearizzazione_MPC.pdf}: modello, linearizzazione,
        costo $J_t$;
  \item {\rifq} $=$ \file{Problema_QP_MPC.pdf}: costruzione delle matrici del
        QP.
\end{itemize}

\paragraph{Come si legge.}
\begin{itemize}
  \item Le voci con il pallino sono \emph{passaggi}, e ciascuna dice se è
        esatta o se introduce un'approssimazione. Le voci con l'etichetta
        \ip{1}, \ip{2}, \dots\ sono \emph{ipotesi}; la
        sezione~\ref{sec:riepilogo} le raccoglie in una tabella.
  \item L'etichetta a destra di una formula dice dove si trova: [L,~(6)] è
        l'equazione (6) di \rif; ``eq.~7'' è l'equazione 7 del paper di Chen,
        Behal e Li (IEEE CDC 2021).
  \item $n = 1,\dots,N$ è la cella, scritta come indice in alto; $k$ è
        l'istante attuale; $j = 1,\dots,p$ il passo lungo l'orizzonte; $\Tc$
        la durata di un passo. Le correnti sono positive in scarica.
  \item La notazione è quella di {\rif} e {\rifq}.
\end{itemize}

% ---------------------------------------------------------------------------
\section{Modello della cella}\label{sec:modello}

\begin{itemize}
  \item[\ip{1}] \emph{Circuito equivalente.} Ogni cella è un generatore
        $V_{oc}$ in serie a una resistenza $R_o$ e a un ramo formato da $R_p$
        e $C_p$ in parallelo. Lo stato è $\zeta = (s, V_p)$: stato di carica
        (SOC) e tensione di polarizzazione. L'uscita è la tensione ai morsetti
        $y$.
  \item[\ip{2}] \emph{Parametri.} $V_{oc}$, $R_o$, $R_p$, $C_p$ dipendono solo
        dal SOC: non dalla temperatura, né dal verso della corrente, né
        dall'invecchiamento. Ciascuno è una tabella nominale su 7 punti di SOC
        $\stab_1 < \dots < \stab_7$, interpolata linearmente. La cella $n$
        differisce dalla nominale per un fattore moltiplicativo $\varphi^n_P$ su
        ogni parametro $P$ e per la capacità $C^n$. Per
        $\stab_\ell \le s < \stab_{\ell+1}$:
        \[
          P^n(s) = \varphi^n_P\,\bigl[P^{\nom}_\ell + m_\ell\,(s - \stab_\ell)\bigr], \quad
          (P^n)'(s) = \varphi^n_P\, m_\ell, \quad
          m_\ell = \frac{P^{\nom}_{\ell+1} - P^{\nom}_\ell}{\stab_{\ell+1} - \stab_\ell} .
          \tag*{[L,~(T)]}
        \]
        Fuori dalla tabella si usa il valore dell'estremo, con derivata nulla.
  \item[\ip{3}] \emph{Collegamento.} Le celle sono in serie e percorse dalla
        corrente di pacco $i$; il bilanciatore aggiunge alla cella $n$ la
        corrente $u^n$. La corrente della cella è
        \[
          I^n = i + u^n .
          \tag*{eq.~3}
        \]
  \item \emph{Modello a tempo continuo}, da \ip{1}--\ip{3}: conteggio della
        carica, bilancio di corrente del ramo RC, tensione ai morsetti.
        \[
        \begin{aligned}
          \dot s(t) &= -\frac{\eta}{3600\,C^n}\, I^n(t) , \\
          \dot V_p(t) &= -\frac{V_p(t)}{\tau\bigl(s(t)\bigr)}
                          + \frac{I^n(t)}{C_p\bigl(s(t)\bigr)} ,
                          \qquad \tau = R_p\,C_p , \\
          y(t) &= V_{oc}\bigl(s(t)\bigr) - V_p(t) - R_o\bigl(s(t)\bigr)\, I^n(t) ,
        \end{aligned}
        \tag*{[L,~(C1)--(C3)]; eq.~1}
        \]
        con $\eta$ il rendimento coulombico (1 nel progetto) e $C^n$ la
        capacità in Ah. Il modello non è lineare perché $\tau$, $C_p$,
        $V_{oc}$ e $R_o$ dipendono da $s$.
  \item[\ip{4}] \emph{Discretizzazione.} La corrente è costante su ogni
        passo: nel passo $j$, cioè per $(j-1)\Tc \le t < j\Tc$, vale
        $I^n(t) = I^n_{k+j-1}$. Le equazioni differenziali si integrano con il
        metodo di Eulero in avanti, con il secondo membro valutato a inizio
        passo.
  \item \emph{Modello discreto}, per $j = 1, \dots, p$:
        \[
        \begin{aligned}
          s_{k+j} &= \min\bigl(s_{k+j-1} - \kappa^n I^n_{k+j-1},\; 1\bigr) , \\
          V_{p,k+j} &= V_{p,k+j-1} + \Tc\left(-\frac{V_{p,k+j-1}}{\tau(s_{k+j-1})}
                      + \frac{I^n_{k+j-1}}{C_p(s_{k+j-1})}\right) , \\
          y_{k+j} &= V_{oc}(s_{k+j}) - V_{p,k+j} - R_o(s_{k+j})\, I^n_{k+j} , \\
          \kappa^n &= \frac{\eta\,\Tc}{3600\,C^n} .
        \end{aligned}
        \tag*{[L,~(2)--(5)]; eq.~2}
        \]
        Esatto per il SOC, il cui secondo membro è costante sul passo.
        Approssimato per $V_p$: a ogni passo $V_p$ è moltiplicata per
        $1 - \Tc/\tau$ invece che per $e^{-\Tc/\tau}$. Il modello discreto è
        stabile se $\Tc < 2\tau$ e non oscilla se $\Tc < \tau$; nel progetto
        $\Tc = 20$~s e $\tau$ è tra circa 30 e 106~s. Il minimo, che impedisce
        a una cella piena di caricarsi oltre, non c'è nell'eq.~2a del paper.
  \item[\ip{5}] \emph{Stato misurato.} All'istante $k$ sono noti lo stato
        $(s^n_k, V^n_{p,k})$ di ogni cella e la corrente di pacco $i_k$: sono
        le condizioni iniziali della predizione. Nel modello Simulink si
        campionano dal modello del pacco; non c'è uno stimatore.
  \item[\ip{6}] \emph{Parametri noti.} La predizione della cella $n$ usa i
        parametri della cella $n$: i fattori $\varphi^n_P$ e la capacità $C^n$ sono
        noti al controllore.
\end{itemize}

% ---------------------------------------------------------------------------
\section{La corrente lungo l'orizzonte}\label{sec:orizzonte}

\begin{itemize}
  \item \emph{Caso generale.} Sia $u^n_{k+j\,|\,k}$ la corrente di
        bilanciamento che all'istante $k$ si prevede di applicare alla cella
        $n$ al passo $k+j$, e $i_{k+j\,|\,k}$ la corrente di pacco prevista
        per lo stesso passo. Con $I^n_{k+j} = i_{k+j\,|\,k} + u^n_{k+j\,|\,k}$
        nel modello discreto, la tensione predetta dipende da tutta la
        sequenza:
        \[
          y^n_{k+j} = y^n_{k+j}\bigl(u^n_{k\,|\,k},\, u^n_{k+1\,|\,k},\, \dots,\,
                      u^n_{k+j\,|\,k}\bigr) , \qquad j = 1, \dots, p .
          \tag*{[L,~(U1)]}
        \]
  \item[\ip{7}] \emph{Corrente di bilanciamento costante sulla predizione.} Si
        sceglie una sola corrente per cella, tenuta per tutto l'orizzonte:
        \[
          u^n_{k+j\,|\,k} = u^n_{k\,|\,k} =: u^n_k , \qquad j = 1, \dots, p .
          \tag*{[L,~(U2)]; eq.~6b}
        \]
        Il paper la giustifica con la lentezza delle correnti di
        bilanciamento.
  \item[\ip{8}] \emph{Corrente di pacco costante sulla predizione.} Del carico
        futuro non si ha un'anteprima:
        \[
          i_{k+j\,|\,k} = i_k , \qquad j = 1, \dots, p .
          \tag*{[L,~(U2)]}
        \]
  \item \emph{Conseguenze} (esatte, date \ip{7} e \ip{8}). La corrente della
        cella è la stessa a ogni passo, e la predizione è funzione di un solo
        numero per cella:
        \[
          I^n_{k+j} = i_k + u^n_k =: I^n , \qquad
          y^n_{k+j} = y^n_{k+j}(u^n_k) , \qquad
          s^n_{k+j} = s^n_{k+j}(u^n_k) .
          \tag*{[L,~(1), (U3)]}
        \]
        Le incognite passano da $N(p+1)$ a $N$. Inoltre, per \ip{1}--\ip{3},
        il modello della cella $n$ contiene solo $I^n$: la predizione della
        cella $n$ dipende solo da $u^n_k$, non dalle correnti delle altre
        celle.
  \item \emph{Orizzonte recedente.} Della soluzione si applica solo il primo
        passo; all'istante $k+1$ il problema si risolve di nuovo con le nuove
        misure, e in generale
        $u^n_{k+1} = u^n_{k+1\,|\,k+1} \ne u^n_{k+1\,|\,k}$. \ip{7} e \ip{8}
        riguardano quindi la predizione, non la corrente applicata al pacco,
        che cambia da un passo all'altro.
\end{itemize}

% ---------------------------------------------------------------------------
\section{Il problema di controllo}\label{sec:problema}

\begin{itemize}
  \item \emph{Grandezza bilanciata.} Il paper scrive i costi per una grandezza
        $\sigma$, che può essere la tensione ($\sigma = y$) o il SOC
        ($\sigma = s$). Qui si scrive $\sigma = y$; con $\sigma = s$ cambia
        solo quanto detto alla fine della sezione~\ref{sec:vincoli}.
  \item \emph{Problema di ottimo non lineare.} All'istante $k$, con
        $u_k = (u^1_k, \dots, u^N_k)$ e per ogni $j = 1, \dots, p$ e
        $n = 1, \dots, N$:
        \begin{align*}
          & \min_{u_k \in \mathbb{R}^N} \;\; J(u_k) \tag*{eq.~6a}\\
          & \text{con il modello discreto, \ip{7}, \ip{8} e lo stato iniziale }
            (s^n_k, V^n_{p,k}) , \tag*{eq.~6b--6c}\\
          & u_{\min} \le u^n_k \le u_{\max} , \tag*{eq.~6d}\\
          & y^n_{k+j}(u^n_k) \ge y_{\min} , \tag*{eq.~6e}\\
          & \textstyle\sum_{n=1}^{N} u^n_k = 0 . \tag*{eq.~6f}
        \end{align*}
  \item[\ip{9}] \emph{Bilanciatore ideale.} Il bilanciatore sposta carica tra
        le celle senza perdite e senza scambiarne con l'esterno, da cui la
        somma nulla; le sue correnti sono limitate (nel progetto
        $u_{\min} = -2$~A e $u_{\max} = 2$~A). Sotto $y_{\min}$ ($3{,}2$~V)
        la cella cede, e con lei il pacco in serie.
  \item \emph{I tre costi.}
        \begin{align*}
          J_t(u_k) &= \sum_{j=1}^{p}\sum_{n=1}^{N}
                      \bigl(y^n_{k+j}(u^n_k) - y^0_{k+j}\bigr)^2
                      + \rho \sum_{n=1}^{N} (u^n_k)^2 , \tag*{eq.~7}\\
          J_m(u_k) &= -\sum_{j=1}^{p} \min_n\, y^n_{k+j}(u^n_k)
                      + \rho \sum_{n=1}^{N} (u^n_k)^2 , \tag*{eq.~8}\\
          J_\Delta(u_k) &= \sum_{j=1}^{p} \Bigl(\max_n\, y^n_{k+j}(u^n_k)
                      - \min_n\, y^n_{k+j}(u^n_k)\Bigr)
                      + \rho \sum_{n=1}^{N} (u^n_k)^2 . \tag*{eq.~11}
        \end{align*}
  \item[\ip{10}] \emph{Peso sulle correnti.} La matrice di peso del paper è
        $R = \rho\, I_N$, con $\rho > 0$ uguale per tutte le celle e
        proporzionale all'orizzonte: $\rho = \rho_{\mathrm{rif}}\, p / p_{\mathrm{rif}}$,
        dove $\rho_{\mathrm{rif}}$ è \texttt{par.mpc.Jt}, \texttt{par.mpc.Jm} o
        \texttt{par.mpc.Jdelta} secondo la formulazione.
  \item[\ip{11}] \emph{Cella nominale.} In $J_t$, $y^0_{k+j}$ è la tensione di
        una cella con parametri e capacità nominali ($\varphi_P = 1$), che parte
        dalla media degli stati misurati e subisce la sola corrente di pacco:
        \[
          s^0_k = \frac1N \sum_{n=1}^{N} s^n_k , \qquad
          V^0_{p,k} = \frac1N \sum_{n=1}^{N} V^n_{p,k} , \qquad
          I^0 = i_k ,
        \]
        fatta avanzare con lo stesso modello discreto. Non dipende da $u_k$:
        per il problema è un dato. Nell'eq.~7 il paper scrive
        $\sigma^0_{k+1}$, da leggere $\sigma^0_{k+j}$.
  \item \emph{Perché non è ancora un QP.} Per tre ragioni, ciascuna con il
        suo rimedio: le tensioni predette non sono funzioni lineari di $u_k$
        (sezione~\ref{sec:lin}); il vincolo eq.~6e può non avere soluzione
        (sezione~\ref{sec:morbido}); il minimo e il massimo di $J_m$ e
        $J_\Delta$ non sono derivabili (sezione~\ref{sec:costo}).
\end{itemize}

% ---------------------------------------------------------------------------
\section{Linearizzazione della predizione}\label{sec:lin}

Le celle sono indipendenti nella predizione: dove non serve, l'indice $n$ è
sottinteso.

\begin{itemize}
  \item[\ip{12}] \emph{Punto di linearizzazione} (Remark~3). Si linearizza
        attorno alla mossa precedente, $\bar u^n = u^n_{k-1}$, supponendo che
        la nuova soluzione $u_k$ ne resti vicina.
  \item \emph{Traiettoria di linearizzazione} (esatta). Il modello discreto
        con la corrente $I^n = i_k + \bar u^n$ dà
        \[
          \bar Y_{j,n} = y^n_{k+j}(\bar u^n) , \qquad
          \bar S_{j,n} = s^n_{k+j}(\bar u^n) , \qquad j = 1, \dots, p .
        \]
  \item[\ip{13}] \emph{Derivabilità.} $y^n_{k+j}(\cdot)$ è derivabile in
        $\bar u^n$. Non lo è solo se la traiettoria cade esattamente su un
        punto della tabella o se il SOC satura a 1: lì si usa la derivata del
        tratto a destra, come nella [L,~(T)], e per la cella satura
        $\alpha = 0$.
  \item \emph{Sviluppo al primo ordine} (approssimato, con un errore di
        ordine $(u^n_k - \bar u^n)^2$):
        \[
        \begin{aligned}
          y^n_{k+j}(u^n_k) &\simeq \bar Y_{j,n} + d_{j,n}\,(u^n_k - \bar u^n) , &
          d_{j,n} &= \left.\frac{\partial y^n_{k+j}}{\partial u^n_k}\right|_{\bar u^n} , \\
          s^n_{k+j}(u^n_k) &\simeq \bar S_{j,n} + \alpha_{j,n}\,(u^n_k - \bar u^n) , &
          \alpha_{j,n} &= \left.\frac{\partial s^n_{k+j}}{\partial u^n_k}\right|_{\bar u^n} .
        \end{aligned}
        \tag*{[L,~(6)]}
        \]
  \item \emph{Derivare rispetto alla corrente di cella} (esatto). Poiché
        $I^n = i_k + u^n_k$ e $i_k$ è fissata da \ip{8},
        $\partial/\partial u^n_k = \partial/\partial I^n$. Si definiscono le
        derivate parametriche dello stato
        \[
          \alpha_j = \frac{\partial s_{k+j}}{\partial I} , \qquad
          \beta_j = \frac{\partial V_{p,k+j}}{\partial I} , \qquad
          \alpha_0 = \beta_0 = 0 ,
        \]
        nulle all'inizio perché lo stato misurato \ip{5} non dipende dalla
        corrente futura. Tutte le derivate si calcolano lungo la traiettoria
        di linearizzazione, cioè con $I = i_k + \bar u^n$.
  \item \emph{Derivata dell'equazione di $V_p$} (esatta). Per la regola della
        funzione composta $\partial\tau/\partial I = \tau'\,\alpha_{j-1}$ e
        $\partial C_p/\partial I = C_p'\,\alpha_{j-1}$, con
        $\tau' = R_p'\,C_p + R_p\,C_p'$ e i parametri letti in $s_{k+j-1}$:
        \[
          \beta_j = \beta_{j-1} + \Tc\left(
             -\frac{\beta_{j-1}}{\tau}
             + \frac{V_{p,k+j-1}\,\tau'\,\alpha_{j-1}}{\tau^2}
             + \frac{1}{C_p}
             - \frac{I\,C_p'\,\alpha_{j-1}}{C_p^2}\right) .
          \tag*{[L,~(7)]}
        \]
  \item \emph{Derivata dell'equazione del SOC} (esatta):
        \[
          \alpha_j =
          \begin{cases}
            \alpha_{j-1} - \kappa & \text{se } s_{k+j-1} - \kappa I < 1 , \\
            0 & \text{se la cella satura a } s_{k+j} = 1 .
          \end{cases}
          \tag*{[L,~(8)]}
        \]
        Senza saturazione $\alpha_j = -j\,\kappa$: il SOC è esattamente lineare
        in $u_k$, e la seconda riga della [L,~(6)] è un'uguaglianza.
  \item \emph{Derivata dell'uscita} (esatta), con i parametri letti in
        $s_{k+j}$:
        \[
          d_j = V_{oc}'\,\alpha_j - \beta_j - R_o'\,\alpha_j\, I - R_o .
          \tag*{[L,~(9)]}
        \]
        Il coefficiente di sensibilità del SOC è $\alpha_j$ stesso: $\alpha$,
        $\beta$ e $d$ sono i coefficienti di sensibilità di $s$, $V_p$ e $y$.
        Sono le derivate esatte del modello discreto, e coincidono con Eulero
        applicato alle equazioni di sensibilità a tempo continuo,
        [L,~(C5)--(C8)]: la linearizzazione non aggiunge un secondo errore di
        discretizzazione.
  \item \emph{Risposta libera e risposta forzata} (riscrittura esatta della
        [L,~(6)]):
        \[
          y^n_{k+j}(u^n_k) \simeq Y^\lib_{j,n} + d_{j,n}\,u^n_k , \qquad
          Y^\lib_{j,n} = \bar Y_{j,n} - d_{j,n}\,\bar u^n ,
          \tag*{[L,~(10)]}
        \]
        e allo stesso modo $s^n_{k+j}(u^n_k) \simeq S^\lib_{j,n} + \alpha_{j,n}\,u^n_k$,
        con $S^\lib_{j,n} = \bar S_{j,n} - \alpha_{j,n}\,\bar u^n$. $Y^\lib_{j,n}$ è
        il valore della retta della cella $n$ in $u^n_k = 0$; non va confusa
        con $y^0_{k+j}$, la tensione della cella nominale.
\end{itemize}

% ---------------------------------------------------------------------------
\section{Incognite e forma del QP}\label{sec:forma}

\begin{itemize}
  \item \emph{Incognite.}
        \[
          z = \begin{bmatrix} u_k \\ \varepsilon \\ \varepsilon_y \end{bmatrix} \in \mathbb{R}^{n_z},
          \qquad n_z = N + n_\varepsilon + 1 ,
        \]
        con $u_k \in \mathbb{R}^N$ le correnti, $\varepsilon \in
        \mathbb{R}^{n_\varepsilon}$ le variabili ausiliarie del costo
        ($n_\varepsilon = 0$ per $J_t$, $p$ per $J_m$, $2p$ per $J_\Delta$;
        sezione~\ref{sec:costo}) ed $\varepsilon_y$ la variabile che allenta il vincolo di
        tensione (sezione~\ref{sec:morbido}).
  \item \emph{Forma di \texttt{quadprog}}:
        \[
          \min_{z} \;\; \tfrac12\, z^\top H z + f^\top z
          \qquad \text{con} \qquad A z \le b , \quad A_{eq}\, z = b_{eq} .
        \]
  \item \emph{Regola di lettura} (esatta; {\rifq}, sezione~2.2). Se il costo
        contiene $a\,z_\ell^2$, allora $H_{\ell\ell} = 2a$; se contiene
        $g\,z_\ell$, allora $f_\ell = g$; le costanti si scartano, perché non
        spostano il punto di minimo. Nei costi che seguono non compaiono
        prodotti $z_\ell z_{\ell'}$ con $\ell \ne \ell'$, quindi $H$ è diagonale.
\end{itemize}

% ---------------------------------------------------------------------------
\section{Vincolo di tensione morbido}\label{sec:morbido}

\begin{itemize}
  \item[\ip{14}] \emph{Vincolo morbido.} Il vincolo eq.~6e è allentato da
        un'incognita $\varepsilon_y \ge 0$, pagata nel costo con un peso $W \gg \rho$ (nel
        progetto $W = 10^3$):
        \[
          y^n_{k+j}(u^n_k) + \varepsilon_y \ge y_{\min} , \qquad \varepsilon_y \ge 0 , \qquad
          \text{costo} \;+\; W \varepsilon_y^2 .
          \tag*{eq.~17}
        \]
        Il problema ha così sempre punti ammissibili, e nella soluzione $\varepsilon_y$
        è la violazione peggiore della predizione:
        \[
          \varepsilon_y = \max\Bigl(0,\; \max_{j,n}\, \bigl(y_{\min} - y^n_{k+j}\bigr)\Bigr) .
        \]
        Il vincolo morbido non è del tutto equivalente a quello rigido: quando una tensione si
        fermerebbe esattamente a $y_{\min}$, si accetta una piccola violazione,
        inversamente proporzionale a $W$. Nel modello $\varepsilon_y$ è sempre presente;
        il paper la introduce solo quando il problema rigido non ha soluzione.
  \item \emph{Contributo al costo} (regola della sezione~\ref{sec:forma}):
        $H_{\varepsilon_y\varepsilon_y} = 2W$, dove
        $H_{\varepsilon_y\varepsilon_y}$ è l'ultimo elemento della diagonale.
\end{itemize}

% ---------------------------------------------------------------------------
\section{Il costo}\label{sec:costo}

\begin{itemize}
  \item \emph{Termine comune.} $\rho\, u_k^\top u_k = \rho \sum_n (u^n_k)^2$
        dà $2\rho$ su ogni elemento $H_{nn}$ delle correnti, per tutti e tre i
        costi.
\end{itemize}

\paragraph{$J_t$, inseguimento della cella nominale.}
\begin{itemize}
  \item \emph{Sostituire la retta} (approssimato, per la [L,~(10)]):
        \[
          y^n_{k+j} - y^0_{k+j} \simeq d_{j,n}\,u^n_k + c_{j,n} , \qquad
          c_{j,n} = Y^\lib_{j,n} - y^0_{k+j} ,
        \]
        \[
          J_t(u_k) \simeq \sum_{j=1}^{p}\sum_{n=1}^{N}
                     \bigl(d_{j,n}\,u^n_k + c_{j,n}\bigr)^2
                     + \rho \sum_{n=1}^{N} (u^n_k)^2 .
          \tag*{[L,~(Q1)]}
        \]
  \item \emph{Separare le celle} (esatto: si scambiano due somme finite):
        \[
          J_t(u_k) = \sum_{n=1}^{N} J^n(u^n_k) , \qquad
          J^n(u^n_k) = \sum_{j=1}^{p} \bigl(d_{j,n}\,u^n_k + c_{j,n}\bigr)^2
                       + \rho\,(u^n_k)^2 .
          \tag*{[L,~(Q2)]}
        \]
        Non ci sono prodotti misti $u^n_k u^m_k$: il blocco delle correnti in
        $H$ è diagonale, e le celle sono accoppiate solo da
        $\sum_n u^n_k = 0$.
  \item \emph{Sviluppare il quadrato} (esatto; $d_j$, $c_j$ e $u$ stanno per
        $d_{j,n}$, $c_{j,n}$ e $u^n_k$):
        \[
          J^n(u) = \Bigl(\sum_{j=1}^{p} d_j^2 + \rho\Bigr) u^2
                 + \Bigl(2\sum_{j=1}^{p} d_j\,c_j\Bigr) u
                 + \sum_{j=1}^{p} c_j^2 .
          \tag*{[L,~(Q3)]}
        \]
        $y^0_{k+j}$ sta dentro $c_j$, quindi compare sia nel coefficiente di
        $u$ sia nella costante: si scarta solo la costante.
  \item \emph{Coefficienti del QP} (regola della sezione~\ref{sec:forma}):
        \[
          H_{nn} = 2\sum_{j=1}^{p} d_{j,n}^2 + 2\rho , \qquad
          f_n = 2\sum_{j=1}^{p} d_{j,n}\,\bigl(Y^\lib_{j,n} - y^0_{k+j}\bigr) .
          \tag*{[L,~(11)]}
        \]
\end{itemize}

\paragraph{$J_m$, massimizzazione della cella più bassa.}
\begin{itemize}
  \item \emph{Riformulazione con minoranti} (esatta; eq.~9--10, Remark~2).
        Per ogni passo si introduce un'incognita $\varepsilon_j$:
        \[
          \min_{u_k,\,\varepsilon,\,\varepsilon_y} \; -\sum_{j=1}^{p} \varepsilon_j
             + \rho\, u_k^\top u_k + W \varepsilon_y^2
          \qquad \text{con} \qquad
          \varepsilon_j \le y^n_{k+j}(u^n_k) \quad \text{per ogni } j, n .
        \]
        Fissata $u_k$, il costo spinge $\varepsilon_j$ verso l'alto e i
        vincoli lo fermano a $\min_n y^n_{k+j}$: le due scritture hanno la
        stessa soluzione.
  \item \emph{Coefficienti}: $f_{N+j} = -1$ e $H_{N+j,\,N+j} = 0$; per le
        correnti $H_{nn} = 2\rho$ e $f_n = 0$.
\end{itemize}

\paragraph{$J_\Delta$, minima differenza tra le celle.}
\begin{itemize}
  \item \emph{Riformulazione con minoranti e maggioranti} (esatta;
        eq.~12--13):
        \[
          \min_{u_k,\,\varepsilon,\,\varepsilon_y} \; \sum_{j=1}^{p} \varepsilon_{p+j}
             - \sum_{j=1}^{p} \varepsilon_j + \rho\, u_k^\top u_k + W \varepsilon_y^2
          \qquad \text{con} \qquad
          \varepsilon_j \le y^n_{k+j}(u^n_k) \le \varepsilon_{p+j} .
        \]
        Nella soluzione $\varepsilon_j = \min_n y^n_{k+j}$ e
        $\varepsilon_{p+j} = \max_n y^n_{k+j}$.
  \item \emph{Coefficienti}: $f_{N+j} = -1$, $f_{N+p+j} = +1$ e zeri sulla
        diagonale di $H$ nelle posizioni delle $\varepsilon$; per le correnti
        $H_{nn} = 2\rho$ e $f_n = 0$.
\end{itemize}

% ---------------------------------------------------------------------------
\section{Vincoli lineari}\label{sec:vincoli}

\begin{itemize}
  \item \emph{Sostituire la retta} (approssimato, per la [L,~(10)]) e portare
        le incognite a sinistra. Per $j = 1, \dots, p$ e $n = 1, \dots, N$:
        \begin{align*}
          -d_{j,n}\,u^n_k - \varepsilon_y &\le Y^\lib_{j,n} - y_{\min}
             && \text{tensione minima} \tag*{[L,~(12)]}\\
          \varepsilon_j - d_{j,n}\,u^n_k &\le Y^\lib_{j,n}
             && \text{minoranti, } J_m \text{ e } J_\Delta \tag*{[L,~(13)]}\\
          d_{j,n}\,u^n_k - \varepsilon_{p+j} &\le -Y^\lib_{j,n}
             && \text{maggioranti, } J_\Delta \tag*{[L,~(14)]}\\
          u_{\min} \le u^n_k &\le u_{\max}
             && \text{limiti del bilanciatore} \tag*{[L,~(15)]}\\
          -\varepsilon_y &\le 0 \tag*{[L,~(16)]}\\
          \textstyle\sum_{n=1}^{N} u^n_k &= 0
             && \text{unico vincolo di uguaglianza} \tag*{[L,~(17)]}
        \end{align*}
        Sono le righe di $A z \le b$ e di $A_{eq}\, z = b_{eq}$; il vincolo
        doppio sui limiti dà due righe per cella.
  \item \emph{Se si bilancia il SOC} ($\sigma = s$). Nel costo e nei vincoli
        sulle $\varepsilon$ si usano $S^\lib_{j,n}$, $\alpha_{j,n}$ e $s^0_{k+j}$
        al posto di $Y^\lib_{j,n}$, $d_{j,n}$ e $y^0_{k+j}$; il vincolo di
        tensione minima usa sempre $Y^\lib_{j,n}$ e $d_{j,n}$.
\end{itemize}

% ---------------------------------------------------------------------------
\section{Il QP risolto a ogni passo}\label{sec:qp}

\begin{itemize}
  \item \emph{Problema finale.}
        \[
          \min_{z} \;\; \tfrac12\, z^\top H z + f^\top z
          \qquad \text{con} \qquad A z \le b , \quad A_{eq}\, z = b_{eq} ,
        \]
        \[
          H = \operatorname{diag}\bigl(H_{11}, \dots, H_{NN},\;
              \underbrace{0, \dots, 0}_{n_\varepsilon},\; 2W\bigr) ,
          \qquad
          H_{nn} =
          \begin{cases}
            2\sum_j d_{j,n}^2 + 2\rho & \text{per } J_t , \\[2pt]
            2\rho & \text{per } J_m \text{ e } J_\Delta .
          \end{cases}
        \]
  \item \emph{Proprietà} (esatte; {\rifq}, sezione~10). $H$ è diagonale con
        elementi non negativi, quindi il QP è convesso e ogni minimo locale è
        globale. Per $J_t$ $H$ è definita positiva. Per tutti e tre i costi la
        soluzione esiste ed è unica.
  \item \emph{Punto di partenza.} In $u_k = \bar u$ la retta coincide con la
        traiettoria di linearizzazione, e $\bar u = u_{k-1}$ rispetta già
        limiti e somma nulla. Un punto ammissibile $z_{\mathrm{iniz}}$ è quindi
        \[
          u_k = \bar u , \qquad
          \varepsilon_j = \min_n \bar Y_{j,n} , \qquad
          \varepsilon_{p+j} = \max_n \bar Y_{j,n} , \qquad
          \varepsilon_y = \max\bigl(0,\; y_{\min} - \min_{j,n} \bar Y_{j,n}\bigr) ,
        \]
        con le $\varepsilon$ solo dove il costo le prevede. \texttt{quadprog}
        (algoritmo \texttt{active-set}) parte da $z_{\mathrm{iniz}}$; se non trova una
        soluzione restituisce $z_{\mathrm{iniz}}$, cioè mantiene la mossa precedente.
  \item \emph{Applicazione.} Al pacco si applicano per un passo le prime $N$
        componenti della soluzione, $u_k$; $u_k$ diventa il punto di
        linearizzazione $\bar u$ del passo successivo.
  \item[\ip{15}] \emph{Una linearizzazione per passo.} Non si itera, cioè non
        si linearizza di nuovo attorno alla soluzione per risolvere ancora: si
        assume che $u_k$ sia abbastanza vicina a $\bar u$ perché la retta
        basti. Con $\Tc = 20$~s e $p = 30$, tre linearizzazioni per passo non
        hanno cambiato i risultati, e il QP dà le stesse durate del problema
        non lineare ({\rif}, sezione~8).
\end{itemize}

% ---------------------------------------------------------------------------
\section{Riepilogo delle ipotesi}\label{sec:riepilogo}

\begin{center}
\small
\begin{tabular}{@{}l >{\raggedright\arraybackslash}p{6.4cm}
                   >{\raggedright\arraybackslash}p{4.4cm} l@{}}
\toprule
 & Ipotesi & Dove interviene & Tipo \\
\midrule
H1  & Circuito equivalente con un ramo RC & modello continuo & modello \\
H2  & Parametri funzione del solo SOC, lineari a tratti, con deviazioni
      moltiplicative & tabelle; $\tau'$, $C_p'$, $V_{oc}'$, $R_o'$ & modello \\
H3  & Celle in serie, $I^n = i + u^n$ & corrente di cella & modello \\
H4  & Corrente costante sul passo, Eulero in avanti & modello discreto &
      approssimazione \\
H5  & Stato e corrente di pacco misurati in $k$ & condizioni iniziali;
      \mbox{$\alpha_0 = \beta_0 = 0$} & misura \\
H6  & Parametri di ogni cella noti al controllore & predizione di ogni cella
      & modello \\
H7  & $u^n_{k+j|k} = u^n_k$ & predizione; $N$ incognite & predizione \\
H8  & $i_{k+j|k} = i_k$ & predizione; $\partial/\partial u^n_k =
      \partial/\partial I^n$ & predizione \\
H9  & Bilanciatore senza perdite e con correnti limitate &
      $\sum_n u^n_k = 0$, limiti & modello \\
H10 & $R = \rho\, I_N$ & $H_{nn}$ & progetto \\
H11 & Cella nominale come riferimento & $c_{j,n}$ e $f_n$ di $J_t$ &
      progetto \\
H12 & $\bar u = u_{k-1}$ & $\bar Y$, $\bar S$, $d$, $\alpha$, $z_{\mathrm{iniz}}$ &
      linearizzazione \\
H13 & Predizione derivabile in $\bar u$ & $d$, $\alpha$ & linearizzazione \\
H14 & Vincolo di tensione morbido & $\varepsilon_y$, $H_{\varepsilon_y\varepsilon_y}$, [L,~(12)] &
      rilassamento \\
H15 & Una linearizzazione per passo & ciclo di controllo &
      linearizzazione \\
\bottomrule
\end{tabular}
\end{center}

\noindent
Le ipotesi \ip{1}--\ip{11} definiscono il problema eq.~6 stesso: valgono
allo stesso modo per il QP e per il problema non lineare risolto con
\texttt{nlmpc}. Il QP se ne discosta solo per i passaggi approssimati delle
sezioni~\ref{sec:lin}--\ref{sec:vincoli}, cioè lo sviluppo al primo ordine
(ipotesi \ip{12}, \ip{13} e \ip{15}) e il vincolo morbido (ipotesi
\ip{14}). Tutti gli altri
passaggi dopo la sezione~\ref{sec:problema} sono esatti: separazione delle
celle, sviluppo dei quadrati, scarto delle costanti, riformulazioni con le
$\varepsilon$.

\end{document}
